\documentclass{beamer}
\usetheme{Madrid}
\usecolortheme{default}

\usepackage[utf8]{inputenc}
\usepackage[T1]{fontenc}
\usepackage{listings}
\usepackage{xcolor}

\title{Unified Wildfire Dataloader Pipeline}
\author{Yangshuang Project Update}
\date{\today}

\lstdefinestyle{pycode}{
  language=Python,
  basicstyle=\ttfamily\scriptsize,
  keywordstyle=\color{blue},
  commentstyle=\color{gray},
  stringstyle=\color{teal},
  showstringspaces=false,
  breaklines=true,
  frame=single
}

\begin{document}

\begin{frame}
  \titlepage
\end{frame}

\begin{frame}{Goal}
Build a single dataloader module that:
\begin{itemize}
  \item Reads heterogeneous local datasets (FIRMS, MTBS FOD, ERA5, NOAA, LANDFIRE, MERRA2)
  \item Normalizes them into one spatial-temporal format
  \item Returns model-ready tensors with a simple API
  \item Keeps raw data untouched (read-only source, optional processed cache)
\end{itemize}
\end{frame}

\begin{frame}{Current API}
\textbf{Primary entrypoint:} \texttt{load\_data(...)}
\begin{itemize}
  \item \texttt{data}: selected sources, e.g. \texttt{["MTBS","FIRMS"]}
  \item \texttt{date\_range}: \texttt{("YYYY-MM-DD","YYYY-MM-DD")}
  \item \texttt{bbox}: \texttt{(min lon, min lat, max lon, max lat)}
  \item \texttt{resolution}: target grid resolution in degrees
  \item optional label controls: \texttt{label\_map}, \texttt{label\_rule}
\end{itemize}
\vspace{0.3em}
\textbf{Return object:}
\begin{itemize}
  \item \texttt{x}: feature tensor \texttt{[T, C, H, W]}
  \item \texttt{y}: label tensor \texttt{[T, H, W]} (or empty/zero placeholder if unavailable)
  \item \texttt{meta}: channels and config metadata
\end{itemize}
\end{frame}

\begin{frame}[fragile]{Example 1: MTBS Label Query (Code)}
\begin{lstlisting}[style=pycode]
from dataloader.adapters.mtbs_fod import MTBSFODAdapter
from dataloader.config import LabelConfig, TimeConfig

bbox = (-87.8, 24.0, -79.8, 31.5)
adapter = MTBSFODAdapter(label_map=LabelConfig().incid_type_map)
df = adapter.load_points(
    "/home/yangshuang",
    TimeConfig(start_date="2022-01-01", end_date="2023-12-31")
)
fl = df[
    (df["longitude"].between(bbox[0], bbox[2])) &
    (df["latitude"].between(bbox[1], bbox[3]))
]
print("rows:", len(fl))
print("label_counts:", fl["label"].value_counts().to_dict())
print(fl.head().to_string(index=False))
\end{lstlisting}
\end{frame}

\begin{frame}[fragile]{Example 1: MTBS Label Query (Output)}
\begin{lstlisting}[style=pycode]
rows: 110
label_counts: {2: 85, 1: 23, 0: 2}
 latitude  longitude       date  label
25.405119 -80.716048 2022-01-31      2
25.768211 -80.432921 2022-03-29      1
26.293827 -81.058581 2022-03-28      1
25.407506 -80.480635 2022-04-26      1
26.290969 -80.521617 2022-08-16      1
\end{lstlisting}
\vspace{0.3em}
Encoding used in this run:
\begin{itemize}
  \item 0 = unknown
  \item 1 = wildfire
  \item 2 = prescribed fire
\end{itemize}
\end{frame}

\begin{frame}[fragile]{Example 2: End-to-End \texttt{load\_data} (Code)}
\begin{lstlisting}[style=pycode]
from dataloader import load_data
from dataloader.adapters import firms as firms_mod

# Fast local run: limit FIRMS to Florida CSVs
firms_mod.FIRMSAdapter.CSV_PATTERNS = ["firmsFL14-25/*.csv"]
firms_mod.FIRMSAdapter.JSON_PATTERNS = []

sample = load_data(
    data=["MTBS", "FIRMS"],
    date_range=("2023-01-01", "2023-01-31"),
    bbox=(-81.0, 26.0, -80.0, 27.0),
    resolution=0.1,
    root_dir="/home/yangshuang",
)
print("x shape:", sample.x.shape)
print("y shape:", sample.y.shape)
print("channels:", sample.meta.get("channels"))
print("x nonzero:", int((sample.x != 0).sum()))
print("y nonzero:", int((sample.y != 0).sum()) if sample.y.size else 0)
\end{lstlisting}
\end{frame}

\begin{frame}[fragile]{Example 2: End-to-End \texttt{load\_data} (Output)}
\begin{lstlisting}[style=pycode]
x shape: (29, 2, 11, 10)
y shape: (0, 11, 10)
channels: ['count', 'frp']
x nonzero: 282
y nonzero: 0
\end{lstlisting}
\vspace{0.3em}
Interpretation:
\begin{itemize}
  \item \texttt{x} contains FIRMS features (count, frp)
  \item \texttt{y} is empty here because no MTBS labels matched this specific
        date-range + bbox alignment
\end{itemize}
\end{frame}

\begin{frame}{Status and Next Steps}
\textbf{Implemented}
\begin{itemize}
  \item Unified API and schema
  \item MTBS FOD label encoding
  \item FIRMS feature rasterization
  \item Adapter scaffolding and first-pass loading for ERA5/NOAA/LANDFIRE/MERRA2
\end{itemize}
\textbf{Planned improvements}
\begin{itemize}
  \item More robust multi-source time alignment strategy
  \item Configurable source filtering in API (no monkey patching)
  \item Stronger validation and source-level diagnostics
\end{itemize}
\end{frame}

\begin{frame}{Takeaway}
The dataloader now provides:
\begin{itemize}
  \item A practical single entrypoint for loading wildfire-related data
  \item Reproducible label extraction from MTBS FOD
  \item Model-ready gridded tensors from local heterogeneous sources
\end{itemize}
\vspace{0.5em}
This gives a concrete foundation for training/inference pipelines and future
feature expansion.
\end{frame}

\end{document}
